\begin{center}\begin{tikzpicture}[font=\small,>=stealth,line join=round]
\draw[gray!40,rounded corners=3pt] (0,0) rectangle (14.2,8.8);\node at (7.1,8.200000000000001) {例3：先求代数解，再保留可以代入的输入};
\node[anchor=west] at (.6,6.8) {比较得到 $t<-\frac12$：};\draw[->,gray] (1,6.3)--(13,6.3);\draw[blue,very thick] (1.3,6.3)--(6.5,6.3);\draw[blue,fill=white] (6.5,6.3) circle(2.7pt);\node[below=7pt] at (6.5,6.3) {$-\frac12$};\node[anchor=west] at (.6,5.0) {输入限制 $-3\le t\le2$：};\draw[->,gray] (1,4.5)--(13,4.5);\draw[violet,very thick] (3,4.5)--(10,4.5);\draw[violet,fill=violet] (3,4.5) circle(2.7pt);\node[below=7pt] at (3,4.5) {$-3$};\draw[violet,fill=violet] (10,4.5) circle(2.7pt);\node[below=7pt] at (10,4.5) {$2$};\node[anchor=west] at (.6,3.1999999999999997) {共同部分 $[-3,-\frac12)$：};\draw[->,gray] (1,2.6999999999999997)--(13,2.6999999999999997);\draw[orange!85!black,very thick] (3,2.6999999999999997)--(6.5,2.6999999999999997);\draw[orange!85!black,fill=orange!85!black] (3,2.6999999999999997) circle(2.7pt);\node[below=7pt] at (3,2.6999999999999997) {$-3$};\draw[orange!85!black,fill=white] (6.5,2.6999999999999997) circle(2.7pt);\node[below=7pt] at (6.5,2.6999999999999997) {$-\frac12$};
\node[text width=12.6cm,align=center] at (7.1,.65) {三行共用同一 $t$ 尺度；最左端可取，右端使和等于0而排除。};
\end{tikzpicture}\end{center}
